\documentclass[11pt,letterpaper]{article}
\usepackage[utf8]{inputenc}
\usepackage{geometry}
\geometry{letterpaper, margin=0.6in}  % Reduced margins

\usepackage{parskip}
\usepackage{titlesec}
\usepackage{enumitem}
\usepackage[colorlinks=true, allcolors=blue]{hyperref}
\usepackage{fontawesome5}
\usepackage{xcolor}
\usepackage{textcomp}
\pagestyle{empty}

% Custom colors
\definecolor{sectionblue}{RGB}{0,60,113}

% Section title formatting - reduced spacing
\titleformat{\section}
{\Large\bfseries\color{sectionblue}}
{}{0em}
{}[\vspace{-0.8em}\rule{\textwidth}{0.4pt}]

\titlespacing{\section}{0pt}{8pt}{4pt}  % Tighter section spacing

% Subsection formatting
\titleformat{\subsection}
{\bfseries}
{}{0em}{}

\titlespacing{\subsection}{0pt}{4pt}{2pt}  % Tighter subsection spacing

% Tighter list spacing
\setlist{nolistsep}
\setlist[itemize]{leftmargin=*,topsep=1pt,parsep=0pt}
\setlist[enumerate]{leftmargin=*,topsep=1pt,parsep=0pt}

\begin{document}

\begin{center}
    {\LARGE\bfseries Xiaoyu Liu} \\[6pt]  % Increased spacing after name
    \color{sectionblue}  % Contact info in blue
    \href{mailto:liu2999@purdue.edu}{\faEnvelope\ liu2999@purdue.edu} 
    \textcolor{black}{\ \textbullet\ } 
    \faPhone\ (765) 715-0648
    \textcolor{black}{\ \textbullet\ }
    \faUniversity\ Purdue University
\end{center}

\section*{Education}
\noindent
% \textbf{Ph.D. in Applied Mathematics} \hfill 2019 -- 2025 (expected)\\
% \textit{Purdue University}, West Lafayette, IN\\
% Joint with Computational Interdisciplinary Graduate Programs\\
% Advisor: Prof. Jianlin Xia\\[6pt]
\textbf{B.S. in Mathematics and Physics}, Minor in Philosophy \hfill 2022 -- 2026\\
\textit{Purdue University}, West Lafayette, IN \\
% Mentor: Prof. Jonathon Peterson, Prof. Xiangxiong Zhang

% \section*{Research Interests}
% \begin{itemize}
%     \item \textbf{Matrix Analysis in Neural Networks}
%     \begin{itemize}
%         \item Development of efficient second-order optimization methods for neural network training through structured matrix analysis
%         \item Design and implementation of fast solvers for structured matrices in deep learning applications
%     \end{itemize}
%     \item \textbf{Randomized Numerical Linear Algebra}
%     \begin{itemize}
%         \item Advanced randomized algorithms for low-rank matrix approximations with emphasis on Nyström methods
%         \item Theoretical analysis and practical implementation of random estimators for matrix properties
%     \end{itemize}
% \end{itemize}

\section*{Publications}
\begin{enumerate}[leftmargin=*]
% \item  Z. Cai, \textbf{T. Ding}, M. Liu, X. Liu, and J. Xia. Fast solvers for shallow R{\scriptsize e}LU neural network approximations: the 2D case \\
% \textit{In preparation for SIAM Journal on Matrix Analysis and Applications} \hfill 2025
%     \item Z. Cai, \textbf{T. Ding}, M. Liu, X. Liu, and J. Xia. ``Matrix Analysis from Shallow R{\scriptsize E}LU Neural Network Least-Squares Approximations'' \\
%     \textit{In preparation for SIAM Journal on Matrix Analysis and Applications} \hfill 2024

% \item Z. Cai, \textbf{T. Ding}, M. Liu, X. Liu, and J. Xia. ``A Structure-Guided Gauss-Newton Method for Shallow R{\scriptsize E}LU Neural Network'' \\
% \textit{Submitted to SIAM Journal on Scientific Computing} \hfill 2024\\
% \href{https://arxiv.org/abs/2404.05064}{arXiv:2404.05064}
\item \textbf{X. Liu} and Z. Wang. ``Convergence of scaled asymptotically-free
self-interacting random walks to Brownian motion perturbed at extrema''\hfill 2024\\
    \textit{To appear on the Electronic Journal of Probability}\\
    \href{https://arxiv.org/abs/2402.11828}{arXiv:2402.11828}

\end{enumerate}

\section*{Research Projects}
\begin{enumerate}[leftmargin=*,label=\arabic*.]
    \item \textbf{One-dimensional branching random walk in random environment}
    \\Advised by Prof. Peterson at Purdue University
    \begin{itemize}
        % \item Study law of large number and higher order fluctuations for position of the rightmost particle
        % \item Try to generalize results from branching simple random walk
        % \item Explore connections to diffusion processes
        \item Investigate branching random walks in a random environment, emphasizing the impact of traps on the maximum particle’s limiting velocity.
        \item Examine high-order corrections of SLLN using large-deviation estimates, environmental tilting and other tools, work toward a central limit theorem.
        \item Explore potential connections to diffusion processes and surface growth models.
        \item Presented in \textit{Purdue Undergraduate Mathematics Research Night}.
    \end{itemize}
    
    % \item \textbf{Top quark pair production at the LHC for quantum experiments}
    % \\Advised by Prof. Jung at Purdue University -- \href{https://www.physics.purdue.edu/jung/}{Jung's group}
    % \begin{itemize}
    %     % \item Implement tests for Bell's inequality using CMS collider data
    %     \item Analyze quantum entanglement in top pair production at the LHC.
    %     \item Implement high-sensitivity tests for Bell inequalities through spin correlation analysis.
    %     % \item Analyze quantum entanglement in top pair production at the LHC using novel observables defined on spin correlations, enabling high sensitivity test for Bell inequalities.
    % \end{itemize}
        \item \textbf{Nonperturbative QFT and the analytical structure of S-matrix}
    \\Advised by Prof. Stamou at TU Dortmund, Germany
    %  -- \href{https://physik.tu-dortmund.de/en/research/research-focus-particle-physics/prof-dr-emmanuel-stamou/}{Prof. Stamou's group at TU Dortmund}
    \begin{itemize}
        % \item Implement tests for Bell's inequality using CMS collider data
        % \item Understanding the analytical structure of S-matrix elements in quantum field theory.
    \item Analyzed large-$z$ behavior in BCFW recursion and its dependence on gauge choice.  
    % \item Verified, through explicit 4-point QCD computations, that only specific factorization channels survive under suitable gauges, consistent with planar large-$N$ decomposition.  
    % \item Investigated the all-line shift formalism (Cohen \textit{et al.}, 2011) linking BCFW recursion to the CSW MHV-vertex expansion.  
    % \item Explored the QED--QCD structural duality via Ozeren--Stirling symmetrization, connecting $N=1$ and $N\!\to\!\infty$ Yang--Mills limits.  
    \item Implemented symbolic verification in \texttt{FORM} and \texttt{MaRTIn} for $2{\to}2$ and $2{\to}3$ scattering amplitudes.  
    \item Explored connections to the Amplituhedron and geometric formulations of scattering amplitudes.

        \item Presented in \textit{Purdue High Energy Physics Seminar}.
    \end{itemize}
\end{enumerate}

\section*{Research Presentations}
\subsection*{Talks}
\begin{itemize}
    \item ``From BCFW Recursions to Amplituhedron''\\
    \textit{Purdue High Energy Physics Seminar} \hfill 2025

    \item ``Branching Random Walk in Random Environment''\\
    \textit{Purdue Undergraduate Mathematics Conference} \hfill 2025

    \item ``Spectral methods in learning theory of boolean functions''\\
    \textit{Purdue Student Probability Seminar} \hfill 2025

    \item ``Convergence of scaled asymptotically-free self-interacting random walks to Brownian motion perturbed at extrema''\\
    \textit{Rose-Hulman Undergraduate Mathematics Conference} \hfill 2024\\
    \textit{Purdue Undergraduate Mathematics Research Night} \hfill 2024

    % \item ``A Structure-Guided Gauss-Newton Method for Shallow R{\scriptsize E}LU Neural Networks''\\
    % \textit{International Conference On Preconditioning Techniques For Scientific and Industrial Applications}, Atlanta \hfill 2024 
    
    % \item ``Randomized Techniques in Numerical Linear Algebra''\\
    % \textit{Association for Women in Mathematics}, Purdue Chapter \hfill 2024
\end{itemize}

\subsection*{Poster Presentations}
\begin{itemize}
    \item ``Convergence of scaled asymptotically-free self-interacting random walks to Brownian motion perturbed at extrema''\\
    \textit{Seminar on Stochastic Processes} \hfill 2025

    \item ``Separating Standard Model and SUSY Dataframes Using Neural Networks''\\
    With R. Patel, C. Suh and Y. Wu \\
    \textit{Purdue Undergraduate Research Exhibition} \hfill 2024

    % \item ``Scaling limits of self-interacting random walks''\\
    % \textit{Purdue Undergraduate Research Exhibition} \hfill 2023
    
    % \item ``Fast Solvers for Neural Network Least-squares Approximations''\\
    % \textit{NSF Computational Mathematics PI Meeting}, Seattle \hfill 2024 
    
    % \item ``Particle Method for the Landau Equation''\\
    % \textit{Purdue Undergraduate Research Exhibition} \hfill 2019
\end{itemize}

% \section*{Teaching Experience}
% \textbf{Teaching Assistant}, Department of Mathematics, Purdue University
% \begin{itemize}
%     \item MA 16600: \textit{Calculus II} -- Led recitation sections and held office hours \hfill Spring 2020
% \end{itemize}

% \textbf{Graduate Course Grader}, Department of Mathematics, Purdue University
% \begin{itemize}
%     \item MA 52700: \textit{Advanced Mathematics for Engineers}
%     \item MA 35100: \textit{Linear Algebra}
%     \item MA 45000: \textit{Abstract Algebra}
% \end{itemize}

\section*{Honors and Awards}
\begin{itemize}
    % \item T.T. Moh Graduate Scholarship Fund, Purdue University \hfill 2019, 2024
    % \item Spira Undergraduate Research Award \hfill 2019
    % \item Jandos Scholarship \hfill 2018
    % \item Jean E. Rubin Memorial Scholarship \hfill 2017, 2018
	\item Jerison Memorial Award in Analysis \hfill  2024
    \item Joel Spira Undergraduate Research Fellowship \hfill  2023
    \item Glen E. Baxter Award with Interests in Probability Theory \hfill  2023
	\item Outstanding Sophomore in Mathematics \hfill  2023
\end{itemize}

\section*{Leadership \& Service}
\textbf{Treasurer}, Society for Industrial and Applied Mathematics (SIAM) Student Chapter \hfill 2024 -- Present
\begin{itemize}
    \item Lead budget planning and financial management for chapter activities
    \item Organize annual SIAM Student Chapter Conference and coordinate logistics 
    % \item Collaborate with Center for Computational and Applied Mathematics (CCAM) to organize research seminars
    \item Foster academic community through professional development and networking events
\end{itemize}

\textbf{Organizer}, Graduate Student Probability Seminar, Purdue University \hfill 2024 -- Present
\begin{itemize}
    \item Spring 2024: Martingale Theory
    \item Fall 2024: Concentration inequalities
    \item Spring 2025: Boolean Functions
    \item Spring 2026: Statistical Mechanics
\end{itemize}
% \begin{itemize}
%     \item Provided academic and cultural guidance to first-year undergraduate students
%     \item Promoted diversity and inclusion in STEM fields through mentorship initiatives
% \end{itemize}

\textbf{Group Leader}, Global Dialogues, Purdue University \hfill 2022 -- Present
\begin{itemize}
    \item Provided academic and cultural guidance to first-year undergraduate students
    \item Promoted diversity and inclusion in STEM fields through mentorship initiatives
\end{itemize}

\section*{Courses and Workshops Attended}
\begin{itemize}
    \item \textbf{IAS Park City Mathematics Institute}: ``Quantum Computation''. Park City, Utah \hfill 2023
    \item \textbf{Midwest Probability Colloquium}, Northwestern University \hfill 2023, 2024, 2025
    \item \textbf{Seminar on Stochastic Processes}, Indiana University, Bloomington \hfill 2025
    \item \textbf{SIAM@Purdue Annual Conference}, Purdue University \hfill 2024, 2025
    \item \textbf{Mathematics Courses at Purdue}: Measure Theory, Probability Theory, Functional Analysis, Partial Differential Equations, Algebraic Topology, Stochastic Differential Equations, Reading in Topological QFT, Large Deviations
    \item \textbf{Physics Courses at Purdue}: Classical Mechanics, Quantum Mechanics, Statistical Mechanics, Research in High Energy Physics, Quantum Field Theory
\end{itemize}

\section*{Technical Skills}
\begin{itemize}
    \item \textbf{Scientific / Symbolic Computing}: MATLAB, Maple, Python (NumPy, SageMath), Lean
    \item \textbf{Languages}: English (Fluent), Mandarin Chinese (Native), German (Basic)
\end{itemize}

\end{document}